\slajdid{02-s012}
\begin{frame}[label=02-s012]{Skraćena funkcionalna tabela}
% element: 02-s012:e04
Po pravilu, funkcionalna tabela sadrži sve moguće vrednosti ulaznih signala, odnosno logičkih promenljivih. Ponekad je možemo napisati u \textbf{skraćenoj formi}, koja i dalje pokriva sve moguće situacije.
\par\vspace{2mm}
% element: 02-s012:e01
Isti sistem sa tri ulaza i jednim izlazom, ali sa drugačijim zahtevom:
\par\vspace{1mm}
\Rezultat{Izlaz je na nivou logičke jedinice ako je \alert{bar jedan ulaz} na nivou logičke jedinice.}
\par\vspace{2mm}
\begin{columns}[c,onlytextwidth]
\begin{column}{.31\textwidth}\centering
% element: 02-s012:e02
\begingroup
\setlength{\tabcolsep}{3pt}
\setlength{\aboverulesep}{1pt}\setlength{\belowrulesep}{1pt}
\renewcommand{\arraystretch}{1.05}
\par\noindent
\begin{tabularx}{\linewidth}{*{4}{>{\centering\arraybackslash}X}}
\toprule C&B&A&F\\\midrule
0&0&0&0\\\cmidrule[.2pt](lr){1-4}
\textcolor{DEcrvena}{b}&\textcolor{DEcrvena}{b}&1&1\\\cmidrule[.2pt](lr){1-4}
\textcolor{DEcrvena}{b}&1&\textcolor{DEcrvena}{b}&1\\\cmidrule[.2pt](lr){1-4}
1&\textcolor{DEcrvena}{b}&\textcolor{DEcrvena}{b}&1\\\bottomrule
\end{tabularx}\par
\endgroup

\end{column}
\begin{column}{.65\textwidth}
% element: 02-s012:e03
\textcolor{DEcrvena}{\textbf{\(b\) — „bilo šta“}}: ta promenljiva može biti 0 ili 1, bez uticaja na izlaz za navedeni red.
\par\vspace{2mm}
Ponekad koristimo i oznaku \(\textcolor{DEcrvena}{X}\), sa istim značenjem.
\end{column}
\end{columns}
\end{frame}
